
% Rosto e sinopse, em uma coluna: a tabela de sinopse tem duas colunas
% de texto e nao cabe em 8,25 cm.
\thispagestyle{empty}

\vspace*{2.2cm}

\begingroup
\sffamily\raggedright
{\fontsize{9}{11}\selectfont\color{tinta2}
Análise de dados abertos do setor elétrico brasileiro \textbullet{} UTFPR\par}
\vspace{1.4em}
{\fontsize{28}{32}\selectfont\bfseries
Relatório\par}
\vspace{0.6em}
{\fontsize{14}{18}\selectfont\color{tinta2}
O corte de geração fotovoltaica no Sistema Interligado Nacional\par}
\vspace{2.2em}
{\fontsize{10}{13}\selectfont
\pe[autor]{nome completo}\par
\vspace{0.3em}
{\color{tinta2}\texttt{opensource.leonardi@gmail.com}}\par
\vspace{0.3em}
{\color{tinta2}setembro de 2026}\par}
\endgroup

\vspace{2.4cm}

\begin{minipage}{15cm}
\setlength{\parskip}{0.6em}
\noindent
Doze análises medem o corte de geração fotovoltaica no Sistema Interligado
Nacional em três eixos: a magnitude do corte, a composição declarada da sua
causa e a sua distribuição entre usinas. Todas usam dados abertos e sem
autenticação, do Operador Nacional do Sistema Elétrico, da ANEEL, do catálogo
público Sentinel-2 e da NASA POWER.

O objetivo inicial era outro. A proposta original detectaria instalações
fotovoltaicas por imagem de satélite para verificar programas de crédito rural.
Duas medidas a redefiniram. A primeira estabelece que 21,6\% da energia
disponível no parque com apuração é cortada, o que separa a geração verificada
do recurso solar. A segunda estabelece que 91\% da micro e minigeração
distribuída ocupa menos de um pixel do Sentinel-2, o que fecha a rota de
verificação por imagem gratuita.

Cinco das doze análises examinam resultados das anteriores. Duas revisam uma
conclusão e três delimitam o escopo de outra. As conclusões revisadas
permanecem no relatório, com a correção correspondente registrada na seção que
a produziu.
\end{minipage}

\vspace{1.6cm}

\begin{minipage}{15cm}
\tabfontw
\begin{tabular}{@{}>{\rag\arraybackslash}p{9.0cm}r@{}}
\toprule
\cab{O que este relatório cobre} & \cab{quantidade} \\
\midrule
Análises com figura e dados-fonte exportados   & 12 \\
Análises que revisam ou delimitam uma anterior &  5 \\
Testes estatísticos que vão a arquivo          & 49 \\
Testes que sobrevivem à correção de Holm       & 16 \\
Conclusões retiradas por análise posterior     &  8 \\
Bases de dados abertas consultadas             &  6 \\
\bottomrule
\end{tabular}
\notas{A correção de multiplicidade é aplicada dentro de cada análise, nunca
entre análises. O apêndice~\ref{ap:inventario} traz janela, amostra, unidade de
repetição e teste declarado de cada uma.}
\end{minipage}

\clearpage
